% Modernised reading — fonds Alexandre Grothendieck, shelfmark 53, pages 1-33
% (the whole folder).
%
% This is an INTERPRETATION, not a transcription. It re-reads the pages in
% today's notation and vocabulary, states the hypotheses the manuscript leaves
% implicit, and drops the critical apparatus so that the mathematics reads
% straight through. Everything the brackets and underlines carried moves into
% a footnote. It is held to being mathematically correct as it stands; where
% the manuscript is loose or mistaken, the true statement is what appears here
% and the footnote says what the page has. In any dispute of substance the
% transcription governs: whoever cites Grothendieck must cite batch-01.fr.tex
% and batch-02.fr.tex.
%
% Pass: Opus 5.5 (claude-opus-5-5), 2026-10-03 — first pass, unchecked against the pages by a human.
% The folder was read whole — two batches, pages 1-33, both transcribed —
% before any of this was written. Page 1 is the folder cover; pages 16-20 are
% a letter to him in another hand (June 1973), skipped by the transcription
% and described here only by its mathematical objection, its writer unnamed.
%
% What the folder is. Notes for a course on algebraic groups: pages 2-15 over
% a field (homogeneity, geometric properties of schemes over a field, pi_0,
% Cartier duality), pages 21-33 over a ring (distributions, comodules,
% finiteness, linear embedding of affine groups and monoids, the functors
% V(M) and W(M)). The letter objects to the course's embedding of affine
% monoids in GL(n); page 21 opens with the rectification (M(n)), and on
% page 30 the M is written over a GL: the same correction.
%
% Notation fixed for the whole document (the page's letters in footnotes):
%  - k a field on pages 2-15, a ring on pages 21-33; k_s a separable closure,
%    \bar k an algebraic closure; pi = Gal(k_s/k), as on the page;
%  - a : G x X -> X the action (the page's pi on pp. 2-3), mu the
%    multiplication of G (his pi on p. 31);
%  - A the affine algebra, Delta : A -> A (x) A its comultiplication (his pi
%    on pp. 15, 24-30), epsilon the counit, sigma the antipode, rho_M : M ->
%    M (x) A a coaction (his u_0, omega_M, u_M, Delta_M);
%  - \check A = Hom_k(A, k) = Dist(G)(k), delta the coproduct of Dist(X)
%    (his Delta on p. 21, delta on p. 22);
%  - V(M) : k' -> Hom_k(M, k') and W(M) : k' -> M (x)_k k' throughout (on
%    p. 25 he calls the first W; pp. 32-33 fix the use kept here);
%  - submodules M', quotients N (his M'', N', M' on pp. 27-29).
%
% Substantive departures from the pages (footnoted where they occur):
%  p. 2, the connected components = irreducible components claim is stated
%  for pseudo-torsors under a group locally of finite type; p. 5, the
%  equivalences on a homomorphism u are stated for G, H locally of finite
%  type with c) on a maximal point of H (the page writes G), and the converse
%  half of the lemma for W open and closed; p. 6, « X régulier » supplied in
%  the first Prop., the perfection hypothesis dropped from the second; p. 7,
%  « X réduit » supplied for density of the smooth locus; the margin question
%  « G réduit => G séparable ? » answered (no, over an imperfect field); p. 9,
%  « X irréductible » supplied in the rectified criterion; p. 12, two
%  implications written backwards (irreducible / geometrically irreducible),
%  « Y connexe » supplied; p. 13, U V^{-1} contains G^0 only when U and V
%  meet G^0 (or a common coset); p. 14, a point of T with values in S' gives
%  a clopen part of X_{S'}, not every clopen part is one; p. 21, epsilon(x) = 1
%  supplied beside Delta(x) = x (x) x; p. 30, the subcoalgebra corollary (4)
%  is the one the proof needs, not Cor. 2.
% Modern names given and footnoted as ours: pseudo-torsor, geometrically
% reduced/irreducible/connected/integral, Jacobian criterion, Galois category,
% scheme of connected components pi_0(X/k), pro-étale, restricted Lie algebra,
% Chevalley-Barsotti-Rosenlicht, Cartier duality, comodule (a right comodule
% today), fundamental theorem of coalgebras, pseudocompact algebra, Hopf
% subalgebra, Ax-Grothendieck. « Dist(G) » here is the whole dual A^*, not
% the distributions supported at e that the name designates today.
% Announced and not established: the struck theorem on the openness of the
% geometrically regular locus (p. 6), the proof of b) => a) on p. 7 (« Esq. »),
% the plans of pp. 4, 5 and 15 (structure of radicial groups through p-Lie
% algebras, Chevalley's extension, the role of Ext^1(G, G_m)), the inclusion
% of G_a in Cartier duality (p. 15, « il faut… »), and the last line of p. 33.
\documentclass[11pt,a4paper]{article}
\input{../preamble/grothendieck}

\folder{53}
\pages{1}{33}
\dating{1973}
\watermark{Édition de démonstration\\interprétation personnelle de l'œuvre}
\foldertitle{Groupes algébriques : notes manuscrites (s.d.), lettre (1973) — lecture modernisée du dossier entier}

\begin{document}

\section*{Résumé}

\begin{resume}
Ces feuillets sont des notes pour un cours sur les groupes algébriques : les
groupes que l'on définit par des équations polynomiales, comme le groupe des
matrices inversibles, celui des rotations, ou plus simplement les nombres non
nuls sous la multiplication. On peut chercher les solutions de ces équations
dans n'importe quel corps, ou même dans n'importe quel anneau, et c'est ce
qui fait la richesse et la difficulté du sujet.

La première moitié du dossier affronte une difficulté qui vient des corps
non algébriquement clos. L'équation $x^{2} + y^{2} = 0$ n'a qu'une solution
réelle, l'origine, et définit sur les réels un objet d'un seul tenant ; sur
les complexes, elle se scinde en deux droites, $x = iy$ et $x = -iy$. Une
propriété comme « être d'un seul tenant » ou « ne pas avoir de point
multiple » peut donc se perdre quand on agrandit le corps. Grothendieck
dresse la liste de ces propriétés --- réduit, irréductible, connexe, intègre,
régulier --- et de leurs versions « géométriques », celles qui survivent à
toute extension du corps ; pour chacune, il dit sur quels corps la version
ordinaire et la version géométrique coïncident. Il construit aussi l'objet
qui compte les morceaux d'un espace même quand ces morceaux ne deviennent
visibles qu'après extension du corps : un « ensemble » de composantes sur
lequel le groupe de Galois agit.

Pour les groupes, une idée simple fait presque tout le travail : un groupe
est homogène. Une translation envoie n'importe quel point sur n'importe quel
autre, de sorte qu'une propriété locale vraie en un point est vraie partout.
Il suffit donc qu'un groupe soit « bon » en un seul point pour être lisse ; et
ses morceaux d'un seul tenant sont aussi ses morceaux irréductibles. La
première partie se clôt sur la dualité de Cartier, qui apparie les groupes
commutatifs finis deux à deux, comme on apparie un espace vectoriel et son
dual.

La seconde moitié change de point de vue. Au lieu du groupe, on regarde les
formes linéaires sur ses fonctions, que Grothendieck appelle ses
distributions, par analogie avec les mesures de l'analyse. Un point du groupe
se retrouve comme une masse de Dirac : une distribution qui respecte les
produits de fonctions. Faire agir le groupe sur un espace vectoriel, c'est
faire agir ces distributions, ou encore munir l'espace d'une « coaction ».
De là sort un énoncé de finitude : tout vecteur vit dans un sous-espace de
dimension finie stable sous le groupe. Et de là, deux conséquences : tout
groupe affine est limite de groupes de matrices, et tout groupe affine de
type fini est un groupe de matrices.

Entre les deux moitiés est classée une lettre d'un auditeur du cours, de juin
1973. Elle relève une erreur : un monoïde --- une loi associative avec
élément neutre, sans inverses --- ne peut pas se plonger dans les matrices
inversibles, et l'exemple tient en deux éléments, $0$ et $1$ sous la
multiplication. La page qui suit la lettre s'ouvre sur la rectification : il
faut toutes les matrices, et non les seules matrices inversibles. On voit
ici, à petite échelle, un cours qui se corrige à la lecture d'un auditeur.

Pour aller plus loin, les noms à chercher sont ceux des schémas en groupes
sur un corps, des propriétés géométriques des schémas, de la dualité de
Cartier, des cogèbres et comodules, et du plongement des groupes affines
dans $\mathrm{GL}_{n}$.

\keywords{group scheme over a field, pseudo-torsor, neutral component, geometrically reduced scheme, geometrically irreducible scheme, geometrically connected scheme, geometrically integral scheme, geometric regularity, smooth scheme, Jacobian criterion, Galois category, scheme of geometric connected components, pro-étale scheme, Cartier duality, restricted Lie algebra, Chevalley–Barsotti–Rosenlicht theorem, distribution algebra, comodule, fundamental theorem of coalgebras, pseudocompact algebra, Hopf subalgebra, linear algebraic monoid, pro-algebraic group, vector group functor}
\end{resume}

\section*{Le fil du dossier, et les conventions}

Le dossier compte trente-trois pages : une chemise (page 1), quatorze pages
de notes de sa main sur les groupes algébriques sur un corps (pages 2 à 15),
une lettre d'une autre main (pages 16 à 20), et treize pages de notes de sa
main sur les groupes affines sur un anneau (pages 21 à 33). Aucune page de
notes n'est datée ni paginée par lui ; la date de l'inventaire, 1973, est
celle de la lettre.

\begin{enumerate}
\item[1.] \emph{Groupes et homogénéité} (pages 2 à 5) : la multiplication est
universellement ouverte, les ouverts $U \cdot V$, les parties stables par la
composante neutre, les composantes connexes ; un plan du cours (page 4) ;
les homomorphismes, leur image fermée, les monomorphismes (page 5).
\item[2.] \emph{Régularité et lissité} (pages 6 à 8) : régularité
géométrique, critère jacobien, lieu lisse ; un groupe géométriquement réduit
est lisse ; composantes irréductibles et connexes d'un groupe.
\item[3.] \emph{Schémas sur un corps} (pages 9 à 14) : un critère
d'irréductibilité géométrique, rectifié ; schémas étales et ensembles
galoisiens ; le schéma $\pi_{0}(X/k)$ des composantes connexes géométriques,
traité deux fois (pages 9-10 et page 14) ; les quatre propriétés
géométriques (pages 11 à 13).
\item[4.] \emph{Dualité de Cartier} (page 15).
\item[5.] \emph{La lettre} (pages 16 à 20).
\item[6.] \emph{Distributions et comodules} (pages 21 à 29) : l'algèbre des
distributions, l'action d'un monoïde affine comme coaction, les foncteurs
$V$ et $W$, comodules et $\check{A}$-modules topologiques, sous-modules
stables, théorème de finitude.
\item[7.] \emph{Groupes affines et matrices} (pages 30 à 33) : un groupe
affine sur un corps est limite projective de groupes de type fini, un
monoïde affine de type fini se plonge dans $\mathrm{M}_{n}$ ; exemples ;
représentabilité de $W(M)$ et de $V(M)$.
\end{enumerate}

Les deux moitiés ne se citent pas, mais elles se rejoignent en un point :
l'énoncé de la page 30, dont la lettre conteste la partie b), et que la
page 21 rectifie.

\emph{Conventions.} Sur les pages 2 à 15, $k$ est un corps, $k_{s}$ une
clôture séparable, $\bar k$ une clôture algébrique, et
$\pi = \mathrm{Gal}(k_{s}/k)$ comme sur la page. $G$ est un schéma en groupes
sur $k$, $G^{0}$ sa composante neutre, $G_{\mathrm{réd}}$ le sous-schéma
réduit. L'action d'un groupe sur un schéma est notée
$a \colon G \times X \to X$ et la multiplication de $G$ est notée
$\mu$.\footnote{La page écrit $\pi$ pour l'action (pages 2 et 3), pour le
groupe de Galois (pages 9 à 14), pour la comultiplication (pages 15, 24 à 30)
et pour la multiplication du groupe (page 31). On ne garde $\pi$ que pour le
groupe de Galois.}
Sur les pages 21 à 33, $k$ est un anneau, $\mathcal{O} = \mathcal{O}_{k}$ le
foncteur $k' \mapsto k'$ sur les $k$-algèbres, $A$ l'algèbre affine,
$\Delta \colon A \to A \otimes_{k} A$ la comultiplication, $\varepsilon$ la
coünité, $\sigma$ l'antipode, $\rho_{M} \colon M \to M \otimes_{k} A$ une
coaction,\footnote{Il écrit successivement $u_{0}$, $\omega_{M}$, $u_{M}$,
$\Delta_{M}$ pour la coaction, et $\pi$ pour la comultiplication.} et
\[
  V(M) \colon k' \longmapsto \mathrm{Hom}_{k}(M, k'), \qquad
  W(M) \colon k' \longmapsto M \otimes_{k} k' .
\]
Le premier est contravariant en $M$, le second covariant ; c'est l'usage
des pages 32 et 33, que l'on suit partout.\footnote{Page 25, il appelle
$W(M)$ le foncteur $k' \mapsto \mathrm{Hom}_{k}(M, k')$, et la transcription
note qu'il écrit tantôt $V$, tantôt $W$ pour lui. Les pages 23, 32 et 33
réservent $W$ à $k' \mapsto M \otimes_{k} k'$ et $V$ au dual ; on fixe cet
usage pour tout le document.}
On note $\check{A} = \mathrm{Hom}_{k}(A, k) = V(A)(k)$ le dual de $A$, et
$\delta$ le coproduit de $\mathrm{Dist}(X)$, dual de la multiplication de $A$.

Un même mot désigne deux choses dans le dossier : \emph{séparé} (la
diagonale est une immersion fermée) et \emph{séparable} (géométriquement
réduit) ; la page abrège les deux en « sép. », et on les distingue partout.

\emph{Annoncé et non établi} : le théorème biffé de la page 6 sur le lieu de
régularité géométrique ; la démonstration de b) $\Rightarrow$ a) de la page 7,
réduite à « Esquisse » ; les plans des pages 4, 5 et 15 ; l'extension de la
dualité de Cartier à $\mathbf{G}_{a}$ (« il faut… ») ; la dernière ligne de
la page 33.

\pagerange{2}{3}
\section*{Groupes, pseudo-torseurs et ouverts $U \cdot V$ (pages 2, 3 et 5)}

\emph{Tout schéma en groupes sur un corps est séparé.} En effet la section
unité $e \colon \mathrm{Spec}\, k \to G$ est une immersion fermée (un point
rationnel d'un schéma sur un corps est fermé), et la diagonale de $G$ s'en
déduit par le changement de base $(x, y) \mapsto x y^{-1}$.\footnote{La page
énonce le corollaire sans démonstration ; celle-ci est la démonstration
usuelle (SGA~3, VI$_{A}$, 0.3).}

\emph{La multiplication $\mu \colon G \times G \to G$ est universellement
ouverte.} L'automorphisme $(g, h) \mapsto (gh, h)$ de $G \times G$ la
transforme en la première projection, et une projection
$X \times_{k} Y \to X$ est universellement ouverte pour des schémas sur un
corps. Plus généralement, appelons \emph{pseudo-torseur} sous $G$ un
$k$-schéma $X$ muni d'une action $a$ telle que
$(g, x) \mapsto (gx, x)$ soit un isomorphisme $G \times X \to X \times X$ ;
alors $a \colon G \times X \to X$ est universellement ouverte, par le même
argument.

\textbf{Proposition.}\quad Soient $X$ un pseudo-torseur sous $G$, $U \subset G$ et $V \subset X$ deux
ouverts.
\begin{itemize}
\item[a)] $U \cdot V$ est ouvert.
\item[b)] $U \cdot V$ est stable par translation par $G^{0}$, donc par
générisation.
\item[c)] $U \cdot V$ est stable par spécialisation.
\item[c')] Si $U \cdot V$ est quasi-compact (par exemple $U$ et $V$
quasi-compacts), ou si $X$ est localement noethérien, $U \cdot V$ est
fermé.\footnote{La page écrit « $U \times V$ est fermé », là où il s'agit de
$U \cdot V$ ; l'interligne ajoute « et $X$ quasi-séparé (p.\,ex.\ $X$ un
ps.\ torseur) ». C'est l'hypothèse qui fait marcher l'argument : un ouvert
quasi-compact d'un schéma quasi-séparé est rétrocompact, donc constructible,
et une partie constructible stable par spécialisation est fermée.}
\end{itemize}

La démonstration de b), à la page 3, se ramène au cas où un point
$x = u_{0} v_{0}$ de $U \cdot V$ est rationnel, avec $u_{0} \in U(k)$ et
$v_{0} \in V(k)$, et montre que $g x \in U \cdot V$ pour $g \in G^{0}$ : il
suffit de trouver $u$ dans $U \cap g W^{-1}$, où $W$ est l'image réciproque
de $V$ par $h \mapsto h x$. Or $T = G^{0} u_{0}$ est irréductible et contient
$u_{0} \in U$ et $g u_{0} \in g W^{-1}$ ; les deux ouverts $U \cap T$ et
$g W^{-1} \cap T$ de $T$ sont non vides, donc se rencontrent.\footnote{Cet
argument suppose $G^{0}$ géométriquement irréductible, ce qui n'est établi
qu'au corollaire 2 de la page 8 ; dans le dossier, l'ordre des feuillets ne
suit pas l'ordre logique. La page écrit ici « $G_{0}$ » pour la composante
neutre.}

\textbf{Lemme.}\quad Soit $G$ opérant sur $X$, de sorte que les morphismes orbitaux
$g \mapsto g x$ soient générisants (c'est le cas pour un pseudo-torseur).
\begin{itemize}
\item[1)] Une partie $W$ de $X$ stable par $G^{0}$ est stable par
générisation et par spécialisation ; si $X$ est localement noethérien, elle
est donc ouverte et fermée.
\item[2)] Réciproquement, une partie ouverte et fermée de $X$ est stable par
$G^{0}$.
\end{itemize}

Pour 1) : si $y \in W$ est rationnel et $x$ générise $y$, le morphisme
orbital $\varphi \colon G^{0} \to X$, $g \mapsto g y$, est générisant et
$G^{0}$ est stable par générisation dans $G$, donc $x \in \varphi(G^{0})
\subset W$. Le complémentaire de $W$ est stable par $G^{0}$ lui aussi, donc
par générisation, et $W$ est stable par spécialisation. Dans un schéma
localement noethérien, une partie stable dans les deux sens est réunion de
composantes connexes, donc ouverte et fermée.\footnote{La page dit « i.e., si
$X$ loc.\ noeth., est-elle ouverte et fermée » sans justifier ; l'argument
par les composantes irréductibles, localement finies, est le nôtre. Le milieu
de la démonstration est repris plusieurs fois sur la page ; on n'en donne que
ce qui se lit.}
Pour 2), qui est en tête de la page 5 : si $W$ est ouverte et fermée,
$a^{-1}(W) \cap (G^{0} \times W)$ est ouverte et fermée dans
$G^{0} \times W$ et contient $e \times W$ ; chaque fibre
$G^{0}_{k(w)}$ de $G^{0} \times W \to W$ est connexe et rencontre
$e \times W$, donc $a^{-1}(W)$ contient $G^{0} \times W$. Cette moitié
n'utilise pas l'hypothèse sur les morphismes orbitaux.\footnote{La page
commence par « Si $W$ est ouvert et fermé », biffe « ouvert et fermé » et
écrit « stable par spécialisation » ; l'argument qu'elle donne (une partie
ouverte et fermée de $G^{0} \times W$ contenant $e \times W$) demande bien
que $W$, et donc $a^{-1}(W)$, soit ouverte et fermée. On énonce donc la forme
première.}

\textbf{Corollaire.}\quad Soit $X$ un pseudo-torseur sous $G$. Alors $X$ est réunion disjointe de
parties à la fois ouvertes, fermées et quasi-compactes, stables par $G^{0}$ ;
en particulier les composantes connexes de $X$ sont quasi-compactes. Si de
plus $G$ est localement de type fini, les composantes connexes de $X$ sont
ses composantes irréductibles.\footnote{La page énonce le corollaire sous
l'hypothèse que $G \times X \to X \times X$ est ouvert (« p.\,ex.\ $X$ un
ps.\ torseur »), et sa seconde moitié, numérotée « 2. », sans hypothèse de
finitude. Pour une action quelconque, les composantes connexes de $X$ ne sont
évidemment pas irréductibles (le groupe trivial opère sur tout schéma) ; on
s'en tient aux pseudo-torseurs. L'énoncé sur les composantes irréductibles
résulte, après extension du corps, du corollaire 2 de la page 8.}

En effet, pour $U$ un voisinage ouvert affine de $e$ et $V$ un ouvert affine
de $X$, $U \cdot V$ est ouvert, fermé et quasi-compact par a) et c'), et $X$
est réunion de tels ouverts.

\pagerange{4}{4}
\section*{Un plan du cours (page 4)}

La page 4, au crayon, esquisse la suite. Sur un corps parfait, un groupe
algébrique se dévisse en
\[
  e \subset G^{0}_{\mathrm{réd}} \subset G^{0} \subset G ,
\]
où $G^{0}_{\mathrm{réd}}$ est un groupe lisse connexe, $G^{0}/G^{0}_{\mathrm{réd}}$
un groupe radiciel (infinitésimal), et $G/G^{0}$ un groupe étale --- un
groupe discret muni d'une action continue de $\pi$ ; $G$ est quasi-compact
si et seulement si $G/G^{0}$ est fini.\footnote{La page écrit « $G/G^{0}$
cpt » et, sous la dernière inclusion, « loc.\ cpts sép. » ; on lit qu'un
groupe étale, vu comme $\pi$-ensemble discret, est quasi-compact si et
seulement s'il est fini. Le dévissage lui-même demande $k$ parfait, pour que
$G_{\mathrm{réd}}$ soit un sous-groupe (page 13).}
La structure des groupes lisses connexes est « assez bien connue » : variétés
abéliennes, groupes semi-simples, tores, groupes vectoriels
$\mathbf{G}_{a}^{s}$, avec des formes tordues ; celle des groupes radiciels
se ramène, par dévissage, aux groupes de hauteur $1$, décrits par leurs
$p$-algèbres de Lie ; restent les groupes discrets. Dans le cas commutatif,
la liste des briques est : variétés abéliennes, $\mathbf{G}_{a}$,
$\mathbf{G}_{m}$, $\mathbb{Z}/n\mathbb{Z}$, $\mu_{p}$,
$\alpha_{p}$.\footnote{Les $p$-algèbres de Lie sont les algèbres de Lie
restreintes de Jacobson ; l'équivalence entre groupes de hauteur $1$ et
$p$-algèbres de Lie de dimension finie est dans SGA~3, VII$_{A}$, et dans
Demazure--Gabriel. Le mot « types » de la première ligne est lu avec doute.}

\pagerange{5}{5}
\section*{Homomorphismes (page 5)}

Soit $u \colon G \to H$ un homomorphisme de schémas en groupes localement de
type fini sur $k$.

\textbf{Corollaire.}\quad Les conditions suivantes sont équivalentes :
\begin{itemize}
\item[a)] $u$ est universellement générisant ; a') $u$ est générisant en un
point ;
\item[b)] $u(G) \supset H^{0}$ ;
\item[c)] $u(G)$ contient un point maximal de $H$.
\end{itemize}
Sous ces conditions, $u(G)$ est une partie ouverte et fermée de $H$.

L'image d'un homomorphisme de groupes localement de type fini sur un corps
est fermée --- c'est le sous-groupe fermé quotient de $G$ par le noyau. Si
$u$ est générisant en un point, son image contient le point générique d'une
composante irréductible, d'où c) ; si $u(G)$, fermé, contient un point
maximal, il contient la composante correspondante, un translaté de $H^{0}$,
et donc $H^{0}$ puisque c'est un sous-groupe, d'où b) ; et si
$u(G) \supset H^{0}$, $u$ se factorise par l'ouvert $u(G)$, réunion de
translatés de $H^{0}$, sur lequel il est fidèlement plat, donc
universellement générisant.\footnote{La démonstration de la page est plus
courte et suit un autre chemin : « a) $\Leftrightarrow$ a') connu,
a') $\Leftarrow$ c) trivial, b) $\Rightarrow$ c) trivial », puis
a) $\Rightarrow$ b) par la stabilité de $u(G)$ sous $H^{0}$ et le lemme de la
page 3. La condition c) y est écrite « un point maximal $x$ de $G$ tel que
$u^{-1}(x) \neq \emptyset$ », où il faut lire $H$. L'hypothèse « localement
de type fini », que la page ne pose pas, est celle sous laquelle l'image est
fermée et le passage au quotient fidèlement plat (SGA~3, VI$_{B}$). Un
énoncé biffé, juste au-dessus, déduisait la platitude de $u$ du cas où $H$
est réduit ; une note de marge demande un argument pour « universellement ».}

\textbf{Corollaire.}\quad Si $G$ est quasi-compact, l'image de $u$ est fermée.

\textbf{Proposition.}\quad Un monomorphisme $u \colon G \to H$, avec $G$ de type fini, est une immersion
fermée.

C'est le théorème de SGA~3, VI$_{B}$, 1.4.2 ; la page l'énonce sans
démonstration, et une note de marge au crayon interroge la stabilité par
spécialisation.

Suit un plan. D'un côté, les groupes étales et leurs limites : groupes
discrets, profinis, localement compacts totalement discontinus, éventuellement
tordus, équivalents aux groupes topologiques correspondants munis d'une
action continue de $\pi$. De l'autre, les groupes connexes : sur un corps
parfait, un groupe lisse connexe est extension canonique d'une variété
abélienne par un groupe lisse connexe affine, et un groupe quasi-compact est
limite projective de groupes de type fini.\footnote{L'extension canonique est
le théorème de Chevalley (démontré aussi par Barsotti et Rosenlicht), connu
depuis les années 1950 ; la page dit « Extension (can.) d'une VA par un groupe
connexe lisse affine ». La limite projective est démontrée, pour les groupes
affines, à la page 30.}

\pagerange{6}{7}
\section*{Régularité géométrique et lissité (pages 6 et 7)}

\textbf{Proposition.}\quad Soient $k$ un corps parfait, $X$ un $k$-schéma régulier et $k'$ une extension
de $k$ ; on suppose $X$ localement de type fini sur $k$, ou $k'$ de type fini
sur $k$, de sorte que les anneaux locaux de $X_{k'}$ soient noethériens.
Alors $X_{k'}$ est régulier.\footnote{La page omet l'hypothèse « $X$
régulier », sans laquelle l'énoncé est faux ; on la rétablit. Le corollaire
qui suit, pour un anneau local régulier $A$ tel que $A \otimes_{k} k'$ soit
essentiellement de type fini, est la même chose en termes d'anneaux.}

\textbf{Proposition.}\quad Soit $X$ un $k$-schéma localement essentiellement de type fini. Les
conditions suivantes sont équivalentes :
\begin{itemize}
\item[a)] pour tout $k$-schéma régulier $Y$ (à anneaux locaux de
$X \times_{k} Y$ noethériens), $X \times_{k} Y$ est régulier ;
\item[b)] pour toute extension $k'$ de $k$, $X_{k'}$ est régulier ;
\item[c)] $X_{\bar k}$ est régulier ;
\item[d)] pour toute extension finie radicielle $k'$ de $k$, $X_{k'}$ est
régulier.
\end{itemize}

On dit alors que $X$ est \emph{géométriquement régulier}, ou \emph{lisse} si
$X$ est localement de type fini.\footnote{La page pose « $k$ ext.\
parfait » dans les hypothèses de cette proposition ; elle est
vraie sans perfection, et elle n'a d'intérêt que sans elle, puisque sur un
corps parfait les quatre conditions se réduisent à « $X$ régulier »
(proposition précédente). On supprime l'hypothèse.}
La page indique les trois arguments : b) $\Rightarrow$ a) parce qu'un
homomorphisme local plat d'anneaux locaux noethériens, à base et fibre
régulières, a un but régulier ; c) $\Rightarrow$ b) parce que la régularité
descend par homomorphisme local fidèlement plat ; d) $\Rightarrow$ c) par
passage à la limite. Ce dernier pas est le plus profond : c'est le théorème
d'EGA~IV, 6.7.7, selon lequel la régularité géométrique se teste sur les
extensions finies radicielles.

Un schéma géométriquement régulier est régulier ; la réciproque vaut pour
tout $X$ si et seulement si $k$ est parfait --- prendre
$X = \mathrm{Spec}\, k'$, $k'$ extension finie radicielle non triviale, qui
est régulier mais dont $k' \otimes_{k} k'$ n'est pas réduit. La régularité
géométrique est stable par extension du corps et par produit. Elle a aussi
un sens en un point $x$ : $\mathcal{O}_{X,x}$ est géométriquement régulier ;
et si $k(x)$ est séparable sur $k$, la régularité de $\mathcal{O}_{X,x}$
entraîne sa régularité géométrique.\footnote{Ce dernier énoncé est exact :
pour $X$ localement de type fini, si $k(x)$ est séparable sur $k$ et
$\mathcal{O}_{X,x}$ régulier, $X$ est lisse en $x$ (EGA~IV, 17.15.5). Un
théorème biffé, à la suite, annonçait que l'ensemble des points où $X$ est
géométriquement régulier est ouvert, et ne se lit plus.}

\textbf{Proposition (critère jacobien).}\quad Un $k$-schéma $X$ localement de type fini est lisse en $x$ si et seulement
s'il existe un voisinage ouvert $U$ de $x$ et un morphisme
$f \colon U \to \mathbb{A}^{n}_{k}$ étale en $x$, c'est-à-dire de
présentation finie, plat et à fibres étales au voisinage de
$x$.\footnote{La page note l'espace affine $E^{n}_{k}$ et met « étale » entre
guillemets ; le nom de critère jacobien est lu avec doute sur la page.}

Par conséquent le lieu lisse de $X$ est ouvert. Il est dense si et seulement
si, en tout point maximal $x$ de $X$, l'anneau artinien $\mathcal{O}_{X,x}$
est réduit et $k(x)$ est une extension séparable de $k$ --- pour une
extension de type fini, cela revient à être une extension finie séparable
d'une extension transcendante pure. En particulier, si $k$ est parfait et
$X$ réduit, le lieu lisse est un ouvert dense.\footnote{La page écrit « Si
$k$ est parfait, l'ens.\ des pts où $X$ est lisse est un ouvert dense », sans
« $X$ réduit » : pour $X = \mathrm{Spec}\, k[\varepsilon]/(\varepsilon^{2})$
le lieu lisse est vide. On ajoute l'hypothèse, que le corollaire précédent
contient d'ailleurs.}

La normalité géométrique se traite de même, sans condition de finitude : pour
$X$ localement de type fini, le lieu où $X$ est géométriquement normal est
ouvert, et dense si $k$ est parfait et $X$ réduit ; une extension de corps
est géométriquement normale si et seulement si elle est séparable.

\pagerange{7}{8}
\section*{Groupes lisses, composantes d'un groupe (pages 7 et 8)}

\textbf{Proposition.}\quad Soit $G$ un schéma en groupes localement de type fini sur $k$. Les
conditions suivantes sont équivalentes :
\begin{itemize}
\item[a)] $G$ est géométriquement réduit en un point ; a') $G$ est
géométriquement réduit ;
\item[b)] $G$ est géométriquement normal ;
\item[c)] $G$ est lisse sur $k$.
\end{itemize}

On a c) $\Rightarrow$ b) $\Rightarrow$ a') $\Rightarrow$ a). Pour
a) $\Rightarrow$ b) : $G$ géométriquement réduit en un point l'est en un point
maximal $g$, où $\mathcal{O}_{G,g} = k(g)$ est un corps séparable sur $k$,
donc géométriquement normal ; après extension du corps, on peut supposer que
$g$ et un autre point $g'$ sont rationnels, et une translation les échange.
C'est l'homogénéité qui fait tout. La même démonstration vaut pour un
$k$-schéma $X$ sur lequel $G$ opère de façon que $G \times X \to X \times X$
soit surjectif.\footnote{La page énonce la proposition sans l'hypothèse
« localement de type fini » pour a) et b), et c) « si $G$ loc.\ de t.f. » ;
on s'en tient au cas localement de type fini, celui de SGA~3, VI$_{A}$,
1.3.1. Le pas de b) à c) est le critère jacobien, et la page ne donne de
b) $\Rightarrow$ a) qu'une « Esquisse de démonstration », sans contenu.}

Une note de marge demande : « Si $G$ de t.f.\ sur $k$, alors $G$ réduit
$\Rightarrow$ $G$ séparable ? » La réponse est non sur un corps imparfait :
pour $k = \mathbb{F}_{p}(t)$, le sous-groupe de $\mathbf{G}_{a}^{2}$ d'équation
$y^{p} = t x^{p}$ est réduit (le polynôme $y^{p} - t x^{p}$ est irréductible)
mais n'est pas lisse, puisque sur $\bar k$ il devient la droite
$y = t^{1/p} x$ comptée $p$ fois.\footnote{La question est laissée ouverte
sur la page ; la réponse et l'exemple sont les nôtres. Sur un corps parfait
la réponse est oui, par le corollaire 1.}

\textbf{Corollaire 1.}\quad Si $k$ est parfait, $G_{\mathrm{réd}}$ est un sous-groupe fermé de $G$,
géométriquement normal, et lisse si $G$ est localement de type fini.

\textbf{Corollaire 2.}\quad Les composantes irréductibles de $G$ sont ses composantes connexes ; en
particulier $G^{0}$ est irréductible.

On peut supposer $k$ parfait et $G$ réduit, donc normal ; or un schéma
normal localement noethérien a pour composantes connexes ses composantes
irréductibles. Sans finitude, la page note que la composante irréductible
de $e$ est stable par générisation, donc contient le spectre de
$\mathcal{O}_{G,e}$.

\textbf{Corollaire 3.}\quad Pour deux ouverts $U$, $V$ de $G$, $U \cdot V$ est stable par translation par
$G^{0}$ ; si $G$ est localement de type fini, ou si $U$ et $V$ sont
quasi-compacts, $U \cdot V$ est ouvert et fermé.

Si $U$ et $V$ rencontrent $G^{0}$, on a $U V^{-1} \supset G^{0}$ : pour
$g \in G^{0}$, $g V \cap G^{0}$ et $U \cap G^{0}$ sont des ouverts non vides
de $G^{0}$, irréductible, donc se rencontrent. En général, pour $u \in U$ et
$v \in V$, le même argument dans les translatés $u G^{0}$ et $G^{0} v$ montre
que $U \cdot V$ contient $u v G^{0}$. L'ensemble $U \cdot V$ est ouvert comme
image de $U \times V$ par $\mu$, universellement ouverte ; il est stable par
spécialisation, et fermé dès qu'il est rétrocompact, ce qui est le cas s'il
est quasi-compact puisque $G$ est séparé.\footnote{La page est un palimpseste
de plusieurs démonstrations encadrées et biffées ; on suit l'ordre qu'indiquent
ses renvois. Dans le cas localement de type fini, $U \cdot V$ est l'image
réciproque d'une partie de $\pi_{0}(G)$, schéma étale, donc ouvert et
fermé sans hypothèse de quasi-compacité.}

\pagerange{9}{10}
\section*{Irréductibilité géométrique, schémas étales, le schéma $\pi_{0}(X/k)$ (pages 9 et 10)}

\emph{Critère d'irréductibilité géométrique, rectifié.} Soient $X$ un
$k$-schéma irréductible et $f \colon Z \to X$ un morphisme. On suppose
\begin{itemize}
\item[a)] $Z$ géométriquement irréductible ;
\item[b)] $f(Z)$ non contenu dans le lieu des points par lesquels passent
plusieurs composantes irréductibles de $X_{\bar k}$ (par exemple, $f(Z)$
contient un point où $X_{\mathrm{réd}}$ est géométriquement normal).
\end{itemize}
Alors $X$ est géométriquement irréductible.\footnote{La page n'écrit pas
« $X$ irréductible », mais la démonstration l'emploie (« $p(X'_{i}) = X$ ») et
l'énoncé est faux sans elle : deux droites disjointes et un point sur l'une
d'elles. On rétablit l'hypothèse.}

Soient $k'$ une extension finie galoisienne de $k$, $p \colon X' = X_{k'} \to X$,
$f' \colon Z' = Z_{k'} \to X'$, et $X'_{i}$ une composante irréductible de $X'$
contenant $f'(Z')$, qui est irréductible. Les transformées de $X'_{i}$ par
$\mathrm{Gal}(k'/k)$ contiennent $f'(Z')$, qui est stable, donc sont égales à
$X'_{i}$ par b). Ainsi $X'_{i}$ est l'image inverse de la partie fermée
irréductible $p(X'_{i})$, qui est $X$ puisque $p$ est ouvert et fermé ; donc
$X' = X'_{i}$ est irréductible. Une extension radicielle ne change pas la
topologie, d'où la conclusion.

Sans b), l'énoncé tombe : $X = \mathrm{Spec}\, \mathbb{R}[x,y]/(x^{2} + y^{2})$
est irréductible, $Z = \mathrm{Spec}\, \mathbb{R}$ s'y envoie sur l'origine,
mais $X_{\mathbb{C}}$ est la réunion des deux droites $x = \pm i y$, qui se
coupent justement à l'origine.

\emph{Schémas étales.} Le foncteur $X \mapsto E = X(k_{s})$ est une
équivalence entre la catégorie $\mathrm{Et}(k)$ des $k$-schémas étales et
celle des ensembles discrets munis d'une action continue de $\pi$. Dans ce
dictionnaire, $X$ quasi-compact de degré $n$ correspond à $E$ fini d'ordre
$n$ ; $X$ connexe (c'est-à-dire $X = \mathrm{Spec}\, K$, $K$ extension finie
séparable) à une action transitive ; la décomposition en composantes
connexes à la décomposition en orbites ; les limites projectives finies et
les limites inductives aux mêmes limites de $\pi$-ensembles ; les
monomorphismes et les épimorphismes aux injections et aux surjections ; et,
pour $\Gamma$ un groupe fini ordinaire, un $\Gamma$-torseur à un
$\pi$-ensemble $E$ muni d'une action à droite simplement transitive de
$\Gamma$ commutant à celle de $\pi$, ce qui revient à la donnée d'un
homomorphisme continu $\pi \to \Gamma$, à conjugaison près.\footnote{C'est
le prototype de ce que Grothendieck appellera catégorie galoisienne
(SGA~1, V). La page écrit « défini par un hom continu $\pi \to G$ » ; la
précision « à conjugaison près » (une fois choisi un point de $E$,
l'homomorphisme est déterminé) est la nôtre. Il note le groupe fini $G$, que
l'on renomme $\Gamma$ pour ne pas le confondre avec le groupe algébrique.}

\emph{Le schéma des composantes connexes.} Supposons $X_{k_{s}}$ localement
connexe, à composantes connexes définies sur des extensions finies de $k$
(par exemple $X$ quasi-compact, ou localement de type fini). On veut un
$k$-schéma $\pi_{0}(X/k)$ dont les points correspondent aux composantes
connexes de $X$, le corps résiduel du point qui correspond à une composante
$X_{(x)}$ étant la plus grande sous-algèbre étale de
$\Gamma(X_{(x)}, \mathcal{O})$.\footnote{La page dit « la plus grande
sous-algèbre séparable de $\Gamma(X_{(x)}, \mathcal{O}_{X_{(x)}})$
sur $k(x)$ » ; c'est « sur $k$ » qu'il faut lire, $k(x)$ étant
précisément ce que l'on définit.}

L'exemple $X = \mathrm{Spec}\, k_{s}$ montre qu'il faut sortir des schémas
étales : $k_{s}$ est la réunion de ses sous-extensions finies galoisiennes
$k_{i}$, $X$ est la limite des $\mathrm{Spec}\, k_{i}$, et
$X_{k_{s}} = \varprojlim (\pi_{i})_{k_{s}}$ est le schéma « constant » de
valeur l'ensemble profini $\pi$ : $\pi_{0}(X_{k_{s}}) = \pi$ comme espace
topologique. D'où la catégorie des schémas \emph{proétales} sur $k$ (entiers
et séparables), équivalente à celle des $\pi$-espaces profinis et à celle des
$k$-algèbres séparables (réunions de sous-algèbres étales), et celle des
schémas localement proétales, équivalente aux $\pi$-espaces localement
compacts totalement discontinus.\footnote{« Proétale » est son mot ; il ne
s'agit pas du site proétale de Bhatt et Scholze (2015), mais des schémas
spectres d'algèbres ind-étales sur un corps.}

Pour $X$ quasi-compact, on construit $\pi_{0}(X/k) = \mathrm{Spec}(A_{0})$ de
deux façons :
\begin{itemize}
\item[a)] $A_{0}$ est la plus grande sous-algèbre ind-étale de
$\Gamma(X, \mathcal{O}_{X})$, et $X \to \mathrm{Spec}(A_{0})$ est le morphisme
canonique ;
\item[b)] $\bar{A}$ est l'algèbre des fonctions localement constantes sur
$X_{k_{s}}$ à valeurs dans $k_{s}$ (réunion des fonctions localement
constantes sur les $X_{k_{i}}$), sur laquelle $\pi$ agit par transport de
structure, et $A_{0} = \bar{A}^{\pi}$.
\end{itemize}
Alors $\pi_{0}(X/k)$ est proétale, $X \to \pi_{0}(X/k)$ est surjectif à
fibres géométriquement connexes, et $\pi_{0}(X/k)$ est fonctoriel en $X$,
commute aux produits et à l'extension du corps de base.\footnote{La
compatibilité aux produits repose sur le théorème de la page 12 (sur un
corps séparablement clos, un produit de schémas connexes est connexe). La
page note $\underline{\pi}_{0}(X/k)$ ; l'algèbre est notée $A$ sur la page,
on l'appelle $A_{0}$ pour la distinguer de l'algèbre affine de la seconde
moitié.}
Plus généralement, $\pi_{0}(X/k)$ se définit comme schéma localement
proétale dès que tout point de $X$ a un voisinage ouvert et fermé $U$ tel
que $U_{\bar k}$ n'ait qu'un nombre fini de composantes connexes.

\pagerange{11}{13}
\section*{Quatre propriétés géométriques des schémas sur un corps (pages 11 à 13)}

Les quatre propriétés suivent le même patron : une caractérisation des corps
sur lesquels la propriété est stable par produit, puis la version
géométrique de la propriété, avec quatre conditions équivalentes, puis la
comparaison, la stabilité par extension du corps et par produit.

\subsection*{Séparabilité}

Un corps $k$ est parfait si et seulement si un produit tensoriel de
$k$-algèbres réduites est réduit, si et seulement si $k' \otimes_{k} k'$ est
réduit pour toute extension finie $k'$. Pour une $k$-algèbre $A$, les
conditions suivantes sont équivalentes : $A \otimes_{k} k'$ réduit pour toute
extension $k'$ ; $A \otimes_{k} B$ réduit pour toute $k$-algèbre réduite $B$ ;
$A \otimes_{k} \bar k$ réduit ; $A \otimes_{k} k'$ réduit pour toute extension
finie radicielle $k'$. On dit alors que $A$ est \emph{séparable}, ou
\emph{géométriquement réduite} ; pour un schéma, la condition se teste sur
un recouvrement affine. Un schéma géométriquement réduit est réduit ; la
réciproque vaut pour tous les $k$-schémas si et seulement si $k$ est
parfait.\footnote{La transcription porte ici une phrase raturée qui se lit
« tt schéma géom.\ [réduit] sur $k$ est [géom.]\ réduit » ; on rétablit le
sens. Une note de marge recommande de « faire gaffe » au passage à la
limite.}

\subsection*{Irréductibilité}

Un corps $k$ est séparablement clos si et seulement si, pour deux
$k$-algèbres intègres $A$ et $B$, $(A \otimes_{k} B)/\mathrm{Nil}$ est
intègre (de façon équivalente : un produit de $k$-schémas irréductibles est
irréductible), si et seulement si, pour toute extension finie $k'$, le noyau
de la multiplication $k' \otimes_{k} k' \to k'$ est nilpotent. Un $k$-schéma
$X$ est \emph{géométriquement irréductible} s'il satisfait les conditions
équivalentes : $X \times_{k} Y$ irréductible pour tout $Y$ irréductible ;
$X_{k'}$ irréductible pour toute extension $k'$ ; $X_{\bar k}$ irréductible ;
$X_{k'}$ irréductible pour toute extension finie $k'$. Un schéma
géométriquement irréductible est irréductible ; la réciproque vaut pour tous
les schémas si et seulement si $k$ est séparablement clos.\footnote{La page
12 écrit l'inverse, « tout $k$-schéma irréductible est géom.\ irréductible,
réciproque vraie ssi $k$ est sép.\ clos » ; la transcription le signale. On
redresse le sens de l'implication, comme dans les trois autres cas.}

\subsection*{Connexité}

\textbf{Théorème.}\quad Sur un corps séparablement clos $k$, le produit de deux $k$-schémas connexes
non vides est connexe.

La page le démontre pour des schémas localement de type fini (le cas général
demande la descente), en utilisant que $X \times_{k} Y \to Y$ est ouvert et
surjectif, de sorte que la topologie de $Y$ est quotient. D'abord $X_{k'}$
est connexe pour toute extension $k'$, car les fibres de $X_{k'} \to X$ sont
les spectres des $k(x) \otimes_{k} k'$, connexes et non vides ; ensuite les
fibres de $X \times_{k} Y \to Y$, qui sont des $X_{k(y)}$, sont connexes et
non vides au-dessus de $Y$ connexe. Réciproquement, si $X \times_{k} X$ est
connexe pour tout $k$-schéma fini connexe $X$, $k$ est séparablement clos.

Un $k$-schéma $X$ est \emph{géométriquement connexe} s'il satisfait les
conditions équivalentes : $X \times_{k} Y$ connexe pour tout $k$-schéma
connexe $Y$ ; $X_{k'}$ connexe pour toute extension $k'$ ; $X_{k_{s}}$
connexe ; $X_{k'}$ connexe pour toute extension finie galoisienne
$k'$.\footnote{La page écrit, en a), « $\forall$ $k$-schéma $Y$ », sans
« connexe » ; l'énoncé serait faux pour $Y$ non connexe. Elle définit aussi
$\bar k$ comme une clôture séparable dans cette proposition, ce qu'on écrit
$k_{s}$.}
Géométriquement connexe entraîne connexe, avec réciproque si et seulement si
$k$ est séparablement clos ; pour une extension $k'$, $X_{k'}$ est
géométriquement connexe sur $k'$ si et seulement si $X$ l'est sur
$k$.\footnote{La page dit « $X$ géom.\ connexe sur $k'$ », là où il faut
« sur $k$ ».}

\subsection*{Intégrité}

Un corps $k$ est algébriquement clos si et seulement si le produit tensoriel
de deux $k$-algèbres intègres est intègre, si et seulement si
$k' \otimes_{k} k'$ est intègre pour toute extension finie $k'$ (ce qui
n'arrive que pour $k' = k$). Un $k$-schéma est \emph{géométriquement intègre}
si $X \times_{k} Y$ est intègre pour tout $Y$ intègre, ou de façon équivalente
si $X_{k'}$ l'est pour toute extension $k'$, ou si $X_{\bar k}$ l'est ; c'est
la conjonction de géométriquement irréductible et géométriquement réduit.
Géométriquement intègre entraîne intègre, avec réciproque si et seulement si
$k$ est algébriquement clos.

\subsection*{Un critère de connexité géométrique, et la composante neutre}

\textbf{Proposition.}\quad Soient $X$ un $k$-schéma connexe et $u \colon Z \to X$ un morphisme, $Z$
géométriquement connexe et non vide (par exemple un point rationnel). Alors
$X$ est géométriquement connexe.

Soit $k'$ une extension finie galoisienne, de groupe $\pi'$ agissant sur
$X_{k'}$ et $Z_{k'}$ de façon compatible à $u_{k'}$. Soit $U$ une partie
ouverte et fermée non vide de $X_{k'}$ ; quitte à la remplacer par son
complémentaire, $U$ contient $u_{k'}(Z_{k'})$, connexe et stable par
$\pi'$. L'intersection des transformées de $U$ par $\pi'$ est ouverte,
fermée, stable et non vide ; elle provient donc d'une partie ouverte et
fermée non vide de $X$, qui est $X$ tout entier. Donc $U = X_{k'}$, et
$X_{k'}$ est connexe ; on conclut par la dernière condition ci-dessus.

\emph{Exemples.} La composante neutre $G^{0}$ contient le point rationnel
$e$ : elle est donc géométriquement connexe, et sa formation commute à
l'extension du corps. Elle est quasi-compacte : pour $U$ un voisinage ouvert
affine de $e$, $U U^{-1}$ est quasi-compact, ouvert et fermé (page 2), et
contient donc la composante connexe $G^{0}$, qui est fermée.\footnote{La
page écrit « [Cor.\ $G^{0}$ quasi-compact] » sans démonstration et biffe
« Elle est quasi-compacte » plus haut ; l'argument par $U U^{-1}$ est celui
de la page 2, rapporté ici.}
Si $k$ est parfait, $G_{\mathrm{réd}}$ est un sous-groupe fermé de $G$, car
$G_{\mathrm{réd}} \times G_{\mathrm{réd}}$ est réduit ; de même
$G^{0}_{\mathrm{réd}}$. Enfin, pour deux ouverts $U$, $V$ de $G$ qui
rencontrent $G^{0}$, $U V^{-1} \supset G^{0}$, et si $U$ et $V$ contiennent
$e$, $U V \supset G^{0}$.\footnote{La page écrit « Soient $U$, $V$ deux ouverts
de $G$, alors $U.V^{-1} \supset G^{0}$ », sans condition : c'est faux si
$U$ et $V$ sont dans deux translatés distincts de $G^{0}$, car alors
$U V^{-1}$ ne rencontre pas $G^{0}$. Il suffit que $U$ et $V$ rencontrent
un même translaté $g G^{0}$.}

\pagerange{14}{14}
\section*{Le schéma des composantes connexes géométriques, seconde mouture (page 14)}

La page 14 reprend la construction des pages 9 et 10 par la descente, au lieu
des algèbres. On suppose $X_{\bar k}$ localement connexe (par exemple
localement noethérien, ce qui est le cas si $X$ est localement de type fini).
Alors $I = \pi_{0}(X_{\bar k})$ est un ensemble discret, et le morphisme
$X_{\bar k} \to I_{\bar k}$ vers le schéma constant de valeur $I$ est
universel pour les morphismes de $X_{\bar k}$ vers les schémas constants. Le
groupe $\pi$ agit sur $I$ ; si l'action est continue --- c'est le cas si les
composantes connexes de $X_{\bar k}$ sont définies sur des extensions finies
de $k$, par exemple si $X$ est quasi-compact ---, $I_{\bar k}$ descend en un
schéma étale $T$ sur $k$, et le morphisme descend en
\[
  f \colon X \longrightarrow T = \pi_{0}(X/k) ,
\]
surjectif à fibres géométriquement connexes, propriété qui caractérise $T$ ;
c'est aussi le morphisme universel de $X$ vers un $k$-schéma étale. Un point
de $T$ à valeurs dans un $k$-schéma $S'$ définit une partie ouverte et fermée
de $X_{S'}$, la fibre de $X_{S'} \to T_{S'}$ au-dessus de la section
correspondante.\footnote{La page dit que ce point « est » une partie ouverte
et fermée de $X_{S'}$, en mettant le verbe entre guillemets. Ce n'est pas une
bijection : pour $X = T = \mathrm{Spec}\, k$, $T(S')$ a un seul élément alors
que $X_{S'} = S'$ peut avoir beaucoup de parties ouvertes et fermées. On
énonce l'injection. Il écrit $\bar k$, où la descente demande une clôture
séparable ; une extension radicielle ne change rien à la topologie.}

Pour $X$ quasi-compact, ou somme de schémas quasi-compacts, sans hypothèse
de continuité, on trouve un $T$ affine proétale aux mêmes propriétés : c'est
le $\pi_{0}(X/k)$ des pages 9 et 10, dont le foncteur des points est « un peu
plus subtil à décrire ».

\pagerange{15}{15}
\section*{Dualité de Cartier (page 15)}

La page s'ouvre sur une question de plan : faut-il travailler sur une base
$S$ quelconque ou sur un corps ? Les groupes non affines sont admis. Pour un
groupe diagonalisable $D(M)$ ($M$ un groupe abélien), les représentations
linéaires s'écrivent comme des modules gradués par $M$, et les homomorphismes
$D(M) \to \mathbf{G}_{m}$ forment le groupe constant $M_{S}$ : $M_{S}$ est
réflexif vis-à-vis de $\mathbf{G}_{m}$.

Le foncteur $\underline{\mathrm{Hom}}_{\mathrm{gr}}(G, \mathbf{G}_{m})$
est-il représentable ? Oui pour $G$ diagonalisable, non en général, par
exemple pour $G = \mathbf{G}_{a}$.\footnote{En caractéristique nulle, les
homomorphismes $\mathbf{G}_{a} \to \mathbf{G}_{m}$ sur une base $S$ sont les
$x \mapsto \exp(\lambda x)$ avec $\lambda$ nilpotent : le foncteur est le
complété formel de $\mathbf{G}_{a}$, qui n'est pas un schéma. La page écrit
« non (p.\,ex.\ $G = \mathbf{G}_{a}$) » dans l'interligne, au-dessus d'un
« affine lisse de t.f. » biffé.}
Mais si $G$ est fini localement libre sur $S$, d'algèbre $\mathcal{A}$, on a
\[
  \underline{\mathrm{Hom}}_{\mathrm{gr}}(G, \mathbf{G}_{m}) \subset
  \underline{\mathrm{Hom}}_{\mathcal{O}\text{-alg}}(\mathcal{O}_{S}[T, T^{-1}], \mathcal{A})
  \simeq W(\mathcal{A})^{\times},
\]
le sous-foncteur des unités $u$ de $\mathcal{A} \otimes \mathcal{O}_{S'}$
telles que $\Delta(u) = u \otimes u$ et $\varepsilon(u) = 1$ (les éléments
de type groupe), puisqu'un homomorphisme de groupes vers $\mathbf{G}_{m}$ est
un homomorphisme de monoïdes vers les unités de $\mathcal{O}$. Ce foncteur
est représenté par $\mathrm{Spec}(\check{\mathcal{A}})$, où
$\check{\mathcal{A}} = \mathrm{Dist}(G)_{S}$, dual de $\mathcal{A}$, est une
algèbre grâce à la cogèbre $\mathcal{A}$ ; cette algèbre est commutative si
$G$ l'est (sinon, c'est son abélianisée qui représente le foncteur) ; et la structure de cogèbre de
$\check{\mathcal{A}}$, duale de la multiplication de $\mathcal{A}$, provient
de ce que le foncteur représenté est un foncteur en groupes.

\textbf{Théorème.}\quad Tout schéma en groupes commutatif fini localement libre $G$ est réflexif pour
$\mathbf{G}_{m}$ : son dual
$G^{*} = \underline{\mathrm{Hom}}_{\mathrm{gr}}(G, \mathbf{G}_{m})$ est
commutatif fini localement libre, et $G \to G^{**}$ est un isomorphisme.
D'où une involution $G \mapsto G^{*}$ de la catégorie de ces groupes, qui
correspond à la dualité des bigèbres bicommutatives finies localement libres
sur $S$.

Exemples : $\mu_{n}$ et $\mathbb{Z}/n\mathbb{Z}$ sont duaux, $\alpha_{p}$ est
son propre dual, et plus généralement $D(M)$ et $M_{S}$ se correspondent.
La dualité marche plus ou moins pour d'autres classes --- groupes discrets,
groupes diagonalisables, formes tordues ---, mais inclure $\mathbf{G}_{a}$
demande autre chose, que la page ne dit pas. Elle « domine la théorie des
groupes algébriques (ou formels) commutatifs, affines ou non », à condition
de la compléter par $\underline{\mathrm{Ext}}^{1}(G, \mathbf{G}_{m})$.\footnote{La
dualité porte le nom de Cartier (1962). Le rôle de
$\underline{\mathrm{Ext}}^{1}(\,\cdot\,, \mathbf{G}_{m})$ est celui qu'il a
pour une variété abélienne $A$, dont la duale est
$\underline{\mathrm{Ext}}^{1}(A, \mathbf{G}_{m})$ (formule de
Weil--Barsotti) ; c'est le cadre de SGA~7, VII et VIII.}

\pagerange{16}{20}
\section*{La lettre de juin 1973 (pages 16 à 20)}

Les pages 16 à 19 sont une lettre en anglais, datée du 15 juin 1973, d'un
auditeur de son cours sur les groupes algébriques ; la page 20 en est le
dernier feuillet, blanc à l'adresse près. Rien n'y est de sa main, et la
transcription ne la reproduit pas ; on n'en dit ici que l'objection
mathématique.\footnote{Ce qui suit est tiré de la description que la
transcription en donne dans son en-tête ; l'auteur de la lettre n'est pas
nommé.}
Elle porte sur la partie b) de « la dernière proposition du cours » : un
monoïde affine de type fini sur un corps se plongerait comme sous-monoïde
fermé de $\mathrm{GL}(n)$. Or l'antipode de la cogèbre de $\mathrm{GL}(n)$
passerait au quotient et ferait du monoïde un groupe. Le contre-exemple est le
monoïde $\{0, 1\}$ sous la multiplication, d'algèbre $k \times k$ ; la lettre
propose de remplacer $\mathrm{GL}(n)$ par $\mathrm{End}_{k}(M)$.

Cette proposition est celle de la page 30, dont le « $\mathrm{M}$ » est écrit
par-dessus un « $\mathrm{GL}$ », et la page 21 s'ouvre sur la
rectification : il faut $\mathrm{M}_{n}$, et non $\mathrm{GL}_{n}$.

\pagerange{21}{22}
\section*{L'algèbre des distributions (pages 21 et 22)}

\emph{Rectification.} Un monoïde affine de type fini sur un corps $k$ admet
une immersion fermée multiplicative dans $\mathrm{M}_{n,k}$, et non dans
$\mathrm{GL}_{n,k}$. Un sous-monoïde fermé de $\mathrm{GL}_{n}$ est en effet
un groupe ; et plus généralement, un monoïde de présentation finie dont les
translations sont des monomorphismes --- par exemple un monoïde qui se plonge
dans un groupe --- est un groupe.\footnote{La fin de la phrase est
incertaine sur la page (« alors $G$ est un objet gr.-pr. »), et
le mot qui qualifie les translations est illisible. Notre lecture : une
translation $G \to G$ qui est un monomorphisme d'un schéma de présentation
finie dans lui-même est surjective (théorème d'Ax--Grothendieck, EGA~IV,
10.4.11, sous sa forme de surjectivité des endomorphismes radiciels), d'où
l'existence d'inverses. La reconstruction est la nôtre.}

Soit $X = \mathrm{Spec}\, A$ un schéma affine sur l'anneau $k$. Ses
\emph{distributions à valeurs dans} $k'$ sont les formes linéaires sur $A$ :
\[
  \mathrm{Dist}(X)(k') = \mathrm{Hom}_{k}(A, k'), \qquad
  \mathrm{Dist}(X) = V(A) \text{ comme $\mathcal{O}$-Module}, \qquad
  \check{A} = \mathrm{Dist}(X)(k) .
\]
\begin{itemize}
\item[a)] $\mathrm{Dist}(X)$ est covariant en $X$.
\item[b)] L'application bilinéaire
$\mathrm{Dist}(X) \times \mathrm{Dist}(Y) \to \mathrm{Dist}(X \times Y)$
est universelle pour les applications bilinéaires vers un $V(M)$ ; en posant
$V(A) \boxtimes V(B) = V(A \otimes_{k} B)$, elle donne
$\mathrm{Dist}(X) \boxtimes \mathrm{Dist}(Y) \simeq \mathrm{Dist}(X \times Y)$,
et $\boxtimes$ coïncide avec $\otimes_{\mathcal{O}}$ si $A$ ou $B$ est
projectif de type fini.
\item[c)] Les distributions ponctuelles (masses de Dirac)
$\alpha_{X} \colon X \to \mathrm{Dist}(X)$ forment une transformation
naturelle compatible aux produits :
$\alpha_{X \times Y}$ est le composé de $\alpha_{X} \times \alpha_{Y}$ et de
l'application bilinéaire de b).\footnote{Le diagramme de la page étiquette la
diagonale $\pi_{X \times Y}$, là où il s'agit de $\alpha_{X \times Y}$ ; la
transcription le signale.}
\end{itemize}

On reconstruit $X$ à partir de la cogèbre $D = \mathrm{Dist}(X)$, munie du
coproduit $\delta \colon D \to D \boxtimes D$ dual de la multiplication de $A$
--- coassociatif, cocommutatif et coünitaire, de coünité $D \to \mathcal{O}$
duale de l'unité ---, comme l'ensemble des $x \in D(k')$ tels que
\[
  \delta(x) = x \boxtimes x \qquad \text{et} \qquad x(1) = 1 ;
\]
ces conditions expriment que la forme linéaire $A \to k'$ est un
homomorphisme d'algèbres.\footnote{La page encadre la seule condition
$\Delta(x) = x \boxtimes x$, qu'elle écrit avec $\Delta$ (et $\delta$ à la
page 22). Elle est satisfaite par $x = 0$ : il faut y joindre la condition de
coünité, que la page évoque dans le crochet « muni de
$D \xrightarrow{\varepsilon} \mathcal{O}$ » sans l'écrire.}

Si $G$ est muni d'une loi de composition, $\mathrm{Dist}(G)$ hérite de la loi
de convolution, associative (commutative) si et seulement si celle de $G$
l'est ; et $G$ a une unité bilatère si et seulement si $\mathrm{Dist}(G)$ a
une unité $e$ telle que $\delta(e) = e \boxtimes e$. Faire opérer $G$ sur un
module $M$, c'est faire opérer $\mathrm{Dist}(G)$ à gauche, d'une façon qui
étend l'action de $G$.

Dans certains cas $\mathrm{Dist}(X)$ se reconstitue à partir de
$\check{A} = \mathrm{Dist}(X)(k)$ seul : si $A$ est projectif de type fini,
ou somme de modules projectifs de type fini (par exemple si $k$ est un
corps), à condition de regarder $\check{A}$ comme module topologique. En
outre $\mathrm{Dist}(G, k)$ contient toute l'information infinitésimale de
$G$ : l'algèbre de Lie de $G$ s'y plonge, son crochet est le commutateur de
l'algèbre, et en caractéristique $p$ sa puissance $p$-ième est la puissance
dans l'algèbre --- de sorte que l'hyperalgèbre, c'est-à-dire la sous-algèbre
des distributions portées par $e$, se lit dans
$\check{A}$.\footnote{La phrase de la page s'arrête sur « de sorte que
l'hyperalgèbre. ». Attention au nom : aujourd'hui, à la suite de
Demazure--Gabriel et de Jantzen, $\mathrm{Dist}(G)$ désigne précisément
cette hyperalgèbre, les formes linéaires sur $A$ nulles sur une puissance de
l'idéal d'augmentation. Le $\mathrm{Dist}(G)(k)$ de ces pages est le dual
tout entier $\check{A}$, qui la contient.}

Exemples. Pour un groupe fini constant, $\check{A}$ est l'algèbre du groupe.
Pour un groupe diagonalisable $D(M)$, $A = k[M]$ est l'algèbre du groupe $M$,
de multiplication induite par la loi de $M$ et de comultiplication induite
par la diagonale $m \mapsto (m, m)$ ; donc
\[
  \mathrm{Dist}(D(M))(k) = k^{M}
\]
avec la multiplication \emph{produit}, terme à terme, qui ne dépend pas de la
loi de groupe de $M$.\footnote{La page glose : « Combinaison linéaire infinie
de caractères ». On lit : un élément de $k^{M}$ est une somme formelle
$\sum_{m} \xi_{m} e_{m}$ des idempotents $e_{m}$ duaux de la base $M$ de $A$,
qui agissent sur une représentation comme les projecteurs sur les
composantes de poids $m$.}

\pagerange{23}{24}
\section*{Une action est une coaction (pages 23 et 24)}

Soient $G = \mathrm{Spec}\, A$ un monoïde affine sur $k$ et $M$ un $k$-module.
Une action de $G$ sur $M$ est un morphisme de monoïdes
$G \to \underline{\mathrm{End}}_{\mathcal{O}}(W(M))$, soumis à $1 x = x$ et
$g(g' x) = (g g') x$. Une « opération sans axiomes » revient à un
endomorphisme $A$-linéaire de $M \otimes_{k} A$, ou encore à une application
$k$-linéaire
\[
  \rho_{M} \colon M \longrightarrow M \otimes_{k} A .
\]
Un point $g \in G(k')$, c'est-à-dire un homomorphisme $g \colon A \to k'$,
agit sur $M \otimes_{k} k'$ par l'endomorphisme $g_{M}$ déduit par changement
de base, et $g_{M}$ est déterminé par sa restriction
$g_{M}^{\circ} = (\mathrm{id}_{M} \otimes g) \circ \rho_{M} \colon M \to M \otimes_{k} k'$.

L'axiome d'unité, pour $k' = k$ et $g = \varepsilon$, devient
\[
  \text{(1')} \qquad (\mathrm{id}_{M} \otimes \varepsilon) \circ \rho_{M} = \mathrm{id}_{M} .
\]
Pour l'associativité, on prend le couple universel : $k' = A \otimes A$,
$g(\lambda) = \lambda \otimes 1$, $g'(\lambda) = 1 \otimes \lambda$ (les deux
projections $G \times G \to G$), de sorte que $g g'$ est la comultiplication
$\Delta$. D'une part $(g g')_{M}^{\circ} = (\mathrm{id}_{M} \otimes \Delta)
\circ \rho_{M}$ ; d'autre part $g'_{M}$ envoie $x$ sur
$\sum y_{i} \otimes 1 \otimes \lambda_{i}$ si
$\rho_{M}(x) = \sum y_{i} \otimes \lambda_{i}$, et $g_{M}$ envoie
$y \otimes 1 \otimes \lambda$ sur $\rho_{M}(y) \otimes \lambda$, de sorte que
$(g_{M} g'_{M})^{\circ} = (\rho_{M} \otimes \mathrm{id}_{A}) \circ \rho_{M}$.
L'axiome 2) s'exprime donc par la commutativité de

\begin{tikzcd}
  M \arrow[r, "\rho_{M}"] \arrow[d, "\rho_{M}"'] & M \otimes A \arrow[d, "\mathrm{id}_{M} \otimes \Delta"] \\
  M \otimes A \arrow[r, "\rho_{M} \otimes \mathrm{id}_{A}"'] & M \otimes A \otimes A
\end{tikzcd}

C'est la définition d'un comodule sur la cogèbre $A$.\footnote{La page écrit
d'abord la coaction $M \to A \otimes_{k} M$, puis échange les facteurs d'un
trait ; on suit l'ordre $M \otimes A$ des formules. Avec la coaction à droite,
c'est ce qu'on appelle aujourd'hui un comodule \emph{à droite} ; la page 26
dit « comodule à gauche », en pensant à l'action de $G$, qui est à gauche.}

\pagerange{25}{27}
\section*{Les foncteurs $V$ et $W$ ; comodules et $\check{A}$-modules (pages 25 à 27)}

\begin{itemize}
\item[a)] Le foncteur $M \mapsto V(M)$, de $\mathrm{Mod}(k)^{\mathrm{op}}$ vers
les $\mathcal{O}$-Modules, est pleinement fidèle :
$\mathrm{Hom}_{k}(M, N) \simeq \mathrm{Hom}_{\mathcal{O}}(V(N), V(M))$.
\item[b)] L'accouplement $V(M) \otimes_{\mathcal{O}} V(N) \to V(M \otimes_{k} N)$
n'est pas un isomorphisme en général, mais il induit une bijection
\[
  \mathrm{Hom}_{\mathcal{O}}(V(M \otimes_{k} N), V(P)) \xrightarrow{\ \sim\ }
  \mathrm{Hom}_{\mathcal{O}}(V(M) \otimes_{\mathcal{O}} V(N), V(P)) ,
\]
et le premier membre est $\mathrm{Hom}_{k}(P, M \otimes_{k} N)$.
\end{itemize}

Il en résulte un dictionnaire. Une loi de composition $\mathcal{O}$-bilinéaire
sur $V(A)$ est une comultiplication $\Delta \colon A \to A \otimes A$ ; une
unité pour cette loi est une coünité $\varepsilon \colon A \to k$ ; la loi est
associative (commutative) si et seulement si $\Delta$ est coassociative
(cocommutative).\footnote{Les mots qualifiant les lois sont illisibles sur la
page ; la transcription restitue le sens, que l'on suit.}
Un accouplement bilinéaire $V(A) \times V(M) \to V(M)$ est une application
$M \to M \otimes A$ ; il est unitaire si et seulement si (1') commute, et
associatif si et seulement si (2') commute. En marge : « on dit aussi : $M$
comodule », ce qui signifie que $V(A)$ opère à gauche sur $W(M)$.

Soient donc $A$ une cogèbre (coassociative, coünitaire) et $M$ un comodule.
Le dual $\check{A} = V(A)(k)$, algèbre pour la convolution
$\xi * \eta = (\xi \otimes \eta) \circ \Delta$, opère à gauche sur $M$ par
\[
  \xi \cdot x = (\mathrm{id}_{M} \otimes \xi)(\rho_{M}(x)) .
\]
Si $A$ est projectif de type fini, $M \otimes A \simeq \mathrm{Hom}_{k}(\check{A}, M)$,
et une structure de comodule sur $M$ équivaut à une structure de
$\check{A}$-module ; de même une structure de cogèbre sur $A$ équivaut à une
structure d'algèbre sur $\check{A}$.

Si $A$ est somme directe de modules projectifs de type fini (par exemple si
$k$ est un corps), on munit $\check{A}$ de la topologie de la convergence
simple ; alors $M \otimes_{k} A \simeq \mathrm{Hom\,cont}_{k}(\check{A}, M)$
pour $M$ discret, et plus généralement
\[
  \mathrm{Hom}(N, M \otimes_{k} A) \simeq
  \mathrm{Hom\,cont}_{k}\bigl(\check{A}, \mathrm{Hom}_{k}(N, M)\bigr),
\]
$\mathrm{Hom}_{k}(N, M)$ étant muni de la convergence simple. Une structure
de comodule sur $M$ revient donc à une structure de $\check{A}$-module telle
que, pour tout $x \in M$, $\xi \mapsto \xi x$ soit continue,
c'est-à-dire s'annule sur un idéal ouvert. Et une structure de cogèbre sur
$A$ revient à une structure d'algèbre \emph{continue} sur $\check{A}$, par le
calcul suivant, où $P = \varinjlim P_{\alpha}$, $Q = \varinjlim Q_{\beta}$,
$R = \varinjlim R_{\gamma}$ sont des réunions filtrantes de sous-modules
projectifs de type fini et $\check{P} = \varprojlim \check{P}_{\alpha}$,
etc. :
\begin{align*}
  \mathrm{Bil\,cont}_{k}(\check{P}, \check{Q}; \check{R})
  &\simeq \varprojlim_{\gamma} \varinjlim_{\alpha,\beta}
     \mathrm{Hom}_{k}(\check{P}_{\alpha} \otimes \check{Q}_{\beta}, \check{R}_{\gamma}) \\
  &\simeq \varprojlim_{\gamma} \varinjlim_{\alpha,\beta}
     \mathrm{Hom}_{k}(R_{\gamma}, P_{\alpha} \otimes Q_{\beta}) \\
  &\simeq \varprojlim_{\gamma} \mathrm{Hom}_{k}(R_{\gamma}, P \otimes Q)
   \simeq \mathrm{Hom}_{k}(R, P \otimes Q) .
\end{align*}
On est dans ce cas, par exemple, si $A$ est réunion filtrante dénombrable de
sous-modules projectifs de type fini, les inclusions admettant des
rétractions. Mais en général on ne peut pas oublier $V(A)$ et ses opérations
au profit du seul $\check{A}$.\footnote{Une algèbre topologique de la forme
$\check{A}$, limite projective stricte d'algèbres de dimension finie, est ce
qu'on appelle depuis Gabriel (1962) une algèbre pseudocompacte ; l'équivalence
entre cogèbres et algèbres pseudocompactes est la forme qu'il donne à cette
dualité.}

\pagerange{27}{29}
\section*{Sous-modules stables et finitude (pages 27 à 29)}

\emph{Quotients.} Soit $M \to N$ un épimorphisme de $A$-comodules en
puissance ; $V(N) \subset V(M)$, et l'on dit que $N$ est stable si $V(N)$ est
stable par l'action de $V(A)$. Cela revient à l'existence d'un
$\rho_{N} \colon N \to N \otimes A$ rendant le carré évident commutatif ; il est
alors unique et satisfait aux axiomes. Si $A$ est l'algèbre d'un monoïde
affine $G$, cela signifie que $G$ opère sur $N$ de façon compatible, et
l'opération est unique.\footnote{La page note le quotient $N'$, puis $M'$ ; on
l'appelle $N$, et les sous-modules $M'$ (la page : $M''$, puis $M'$ à la
page 29).}

\emph{Sous-modules.} Pour $i \colon M' \hookrightarrow M$, il n'est plus vrai
en général que $V(M) \to V(M')$ soit surjectif, et l'opération « induite »
n'est plus évidemment unique. Il s'agit de trouver
$\rho_{M'} \colon M' \to M' \otimes A$ avec
$(i \otimes \mathrm{id}_{A}) \circ \rho_{M'} = \rho_{M} \circ i$ ; elle est
unique si $i \otimes \mathrm{id}_{A}$ est injectif, et satisfait alors aux
axiomes si de plus $i \otimes \mathrm{id}_{A \otimes A}$ est injectif. C'est
le cas si $A$ est plat (le cas courant), si $M/M'$ est plat, ou si $M'$ est
facteur direct dans $M$. Sous ces conditions, $M'$ est stable si et
seulement si $M/M'$ l'est ; et si $A$ est plat, une intersection finie de
sous-modules stables est stable, comme une somme quelconque.\footnote{La
remarque sur les intersections et les sommes est cerclée sur la page, sans
hypothèse ; pour les intersections, il faut que $\otimes A$ commute aux
intersections finies, ce que donne la platitude de $A$. Les sommes sont
stables sans hypothèse.}

\textbf{Proposition.}\quad Si $A$ est somme de modules projectifs de type fini, un sous-$k$-module d'un
$A$-comodule $M$ est stable si et seulement si c'est un sous-$\check{A}$-module.

En effet, si $(a_{j})$ est une famille génératrice adaptée et $(\xi_{j})$ les
formes coordonnées, $\rho_{M}(x) = \sum_{j} (\xi_{j} \cdot x) \otimes a_{j}$.

\textbf{Corollaire 1.}\quad Tout sous-module $M'$ de $M$ est contenu dans un plus petit sous-module
stable $M'^{S}$, le sous-$\check{A}$-module engendré ; si $M'$ est de type
fini, $M'^{S}$ aussi.\footnote{L'exposant $S$ est une lecture de la lettre
tracée, qui ressemble aussi à un $G$. La finitude vient de ce que
$\check{A} \cdot x$ est contenu dans le module engendré par les $y_{i}$ de
$\rho_{M}(x) = \sum y_{i} \otimes \lambda_{i}$ --- argument que la page
n'écrit pas.}

\textbf{Corollaire 2.}\quad Tout comodule $M$ est réunion filtrante de ses sous-modules de type fini
stables.

\textbf{Corollaire 3.}\quad Regardons $A$ comme comodule à gauche et à droite sur elle-même, par
$\Delta$. Tout sous-module $B$ de $A$ est contenu dans un plus petit
sous-module $B^{S}$ stable des deux côtés (le sous-$\check{A}$-bimodule
engendré), de type fini si $B$ l'est ; donc $A$ est réunion filtrante de
sous-modules de type fini stables des deux côtés.\footnote{Le passage est
remanié à l'extrême : un premier « Corollaire 3 » est biffé, avec plusieurs
amorces, et un ajout interlinéaire parle de « la coalgèbre produit
($A^{\circ}$ l'opposée de $A$) », c'est-à-dire de la cogèbre
$A \otimes A^{\circ}$ dont les comodules sont les bicomodules. La syntaxe
exacte n'est pas restituable ; on donne l'énoncé repris en clair.}

Or la stabilité des deux côtés signifie que $\Delta(B) \subset (B \otimes A)
\cap (A \otimes B)$, et sur un corps cette intersection est $B \otimes B$ :
$B$ est une sous-cogèbre.\footnote{La page invoque la platitude de $A/B$ ;
sur un corps, $(B \otimes A) \cap (A \otimes B) = B \otimes B$ pour tous
sous-espaces, ce qui suffit au corollaire 4.}

\textbf{Corollaire 4.}\quad Toute cogèbre $A$ sur un corps est réunion filtrante de sous-cogèbres de
dimension finie ; donc $\check{A}$ est limite projective filtrante stricte
d'algèbres de dimension finie.

C'est ce qu'on appelle aujourd'hui le théorème fondamental des
cogèbres.\footnote{Il figure, sous ce nom, dans Sweedler, \emph{Hopf
Algebras} (1969) ; il était connu auparavant, et le dossier ne cite
personne.}

\pagerange{30}{31}
\section*{Groupes affines et matrices ; l'algèbre des distributions (pages 30 et 31)}

\textbf{Proposition.}\quad Soit $G$ un groupe (resp.\ un monoïde) affine sur un corps $k$.
\begin{itemize}
\item[a)] $G$ est limite projective filtrante stricte de groupes (resp.\
monoïdes) affines de type fini sur $k$.
\item[b)] Si $G$ est de type fini, il existe un entier $n$ et une immersion
fermée compatible aux multiplications $G \to \mathrm{M}(n)_{k}$ (resp.\
$G \to \mathrm{GL}(n)_{k}$ si $G$ est un groupe).\footnote{Le « $\mathrm{M}$ »
est écrit par-dessus un « $\mathrm{GL}$ » : c'est sans doute la trace de la
correction qu'appelle la lettre des pages 16 à 19, et que la page 21
énonce. Le mot « entier » est lu d'après le sens.}
\end{itemize}

Pour a) : $A$ est réunion filtrante de ses sous-cogèbres de dimension finie
(corollaire 4), et la sous-algèbre engendrée par une sous-cogèbre est une
sous-bigèbre, puisque $\Delta$ est un homomorphisme d'algèbres. Si $G$ est un
groupe, d'antipode $\sigma$, et $M$ une sous-cogèbre, $\sigma M$ en est une
aussi --- $M$ est un sous-$\check{A}$-module à gauche (à droite) si et
seulement si $\sigma M$ en est un à droite (à gauche) --- ; donc
$M + \sigma M$ est une sous-cogèbre stable par $\sigma$, puisque
$\sigma^{2} = \mathrm{id}$ pour $A$ commutative, et la sous-algèbre qu'elle
engendre est une sous-algèbre de Hopf. Ainsi $A$ est réunion filtrante de
sous-algèbres de Hopf $A_{i}$ de type fini, et $G = \varprojlim G_{i}$ avec
$G_{i} = \mathrm{Spec}\, A_{i}$.\footnote{La page renvoie au « corollaire 2 »
pour pouvoir supposer les $A_{i}$ stables par la comultiplication ; c'est le
corollaire 4, sur les sous-cogèbres, qui le donne : une sous-algèbre engendrée
par un sous-comodule d'un seul côté n'est pas une sous-bigèbre. Les mots
« filtrante » et « stricte » sont des ajouts incertains de l'interligne.}

Pour b) : $A$ est de type fini ; on choisit un sous-espace de dimension
finie $M$ qui engendre $A$ comme algèbre, et on peut le supposer stable par
la coaction $\Delta$ (corollaire 2), de sorte que $G$ opère sur $M$ par
translations. D'où un morphisme de monoïdes
$G \to \underline{\mathrm{End}}_{k}(M) \simeq \mathrm{M}(n)_{k}$, à valeurs
dans $\underline{\mathrm{Aut}}_{k}(M) \simeq \mathrm{GL}(n)_{k}$ si $G$ est un
groupe, avec $n = \dim M$. Pour $\varphi \in M$, la fonction coefficient
$g \mapsto \varepsilon(g \cdot \varphi) = (g \varphi)(1) = \varphi(g)$ est
l'image d'une fonction coordonnée de $\underline{\mathrm{End}}_{k}(M)$ ; ainsi
$M$ est contenu dans l'image de l'algèbre affine de
$\underline{\mathrm{End}}_{k}(M)$ dans $A$, qui est donc surjective, et le
morphisme est une immersion fermée.

\emph{L'algèbre des distributions.} Si $A$ est l'algèbre affine d'un monoïde
affine $G$ --- une algèbre commutative munie d'une comultiplication qui est
un homomorphisme d'algèbres --- les éléments de
$V(A)(k') = \mathrm{Hom}_{k}(A, k')$ s'appellent les distributions de $G$ à
coefficients dans $k'$ ; pour la convolution, elles forment une $k'$-algèbre
$\mathrm{Dist}(G)(k')$, « l'algèbre du monoïde $G$ à coefficients dans $k'$ »,
et l'on a un homomorphisme injectif de monoïdes
\[
  G(k') \hookrightarrow \mathrm{Dist}(G)(k') ,
\]
à valeurs dans les éléments inversibles si $G$ est un groupe, et
« pratiquement jamais » bijectif.\footnote{La page note
$\mathrm{Dist}(G)(k')^{\times}$ le monoïde multiplicatif et
$\mathrm{Dist}(G)(k')^{*}$ son groupe des unités. Un exemple, celui d'un
groupe fini constant $\Gamma$ dont l'algèbre affine est $k^{\Gamma}$, est
encadré puis biffé.}
L'antipode s'écrit par le diagramme $g g^{-1} = 1$ : le composé
$G \xrightarrow{(\mathrm{id}, \sigma)} G \times G \xrightarrow{\mu} G$ est
le composé $G \to e \to G$ ; dualement,
$m \circ (\mathrm{id}_{A} \otimes \sigma) \circ \Delta = \eta \circ \varepsilon$,
où $m(x \otimes y) = x y$ et $\eta \colon k \to A$ est l'unité.\footnote{La
page dessine les deux triangles, $G \to G \times G \to G$ par $e$ et
$A \leftarrow A \otimes A \leftarrow A$ par $k$, avec la règle
$x \, \sigma y \leftarrow\!\shortmid (x, y)$ ; elle note $\pi$ la
multiplication de $G$ et $p$ l'unité de $A$.}

\pagerange{32}{33}
\section*{Exemples ; les foncteurs vectoriels $W(M)$ et $V(M)$ (pages 32 et 33)}

\emph{Exemples de groupes algébriques affines} : $\mathbf{G}_{m}$,
$\mathbf{G}_{a}$, $\mathrm{GL}(n)$ et ses sous-groupes $\mathrm{SL}(n)$, le
groupe orthogonal $O(Q)$, et plus généralement les stabilisateurs de
tenseurs ; le groupe $\mathrm{Tr}(n)$ des matrices triangulaires, son
sous-groupe $\mathrm{Tr}_{0}(n)$ ; le groupe diagonal $\mathrm{Diag}(n)$ et
son sous-groupe $\mathrm{Diag}_{0}(n) \simeq \mathbf{G}_{m}$ ; l'immersion
de $\mathbf{G}_{a}$ dans $\mathrm{GL}(2)$, $x \mapsto
\bigl(\begin{smallmatrix} 1 & x \\ 0 & 1 \end{smallmatrix}\bigr)$.\footnote{La
page ne définit ni $\mathrm{Tr}_{0}(n)$ ni $\mathrm{Diag}_{0}(n)$ ; on lit
$\mathrm{Tr}_{0}(n)$ comme le sous-groupe unipotent (diagonale égale à $1$)
et $\mathrm{Diag}_{0}(n)$ comme les matrices scalaires, d'après
l'isomorphisme avec $\mathbf{G}_{m}$ qu'elle dessine. La matrice de
l'immersion de $\mathbf{G}_{a}$ est la nôtre ; la page n'écrit que les mots.}

\textbf{Proposition.}\quad Le foncteur $W(M) \colon k' \mapsto M \otimes_{k} k'$, faisceau de
$\mathcal{O}$-Modules, est représentable si et seulement si $M$ est
projectif de type fini.

S'il est représentable, il transforme injections d'algèbres en injections,
donc $M$ est plat ; et il transforme produits d'algèbres en produits, d'où
$M \otimes_{k} k^{I} \simeq M^{I}$ pour tout ensemble $I$, ce qui entraîne
que $M$ est de type fini. Écrivant $0 \to R \to L \to M \to 0$ avec $L$ libre
de type fini, la platitude de $M$ garde la suite exacte après
$\otimes_{k} k^{I}$, d'où $R \otimes k^{I} \simeq R^{I}$ et $R$ de type fini ;
$M$ est donc plat de présentation finie, c'est-à-dire projectif de type
fini.\footnote{Le passage de la préservation des injections à la platitude
n'est pas détaillé sur la page ; il suffit, pour un idéal $J$ de $k$, de
plonger l'algèbre $k \oplus J$ (de carré nul sur $J$) dans
$k[\varepsilon] = k \oplus k \varepsilon$, ce qui donne l'injectivité de
$M \otimes J \to M$.}
Inversement, si $M$ est facteur direct de $L = k^{n}$, image d'un
projecteur $p$, alors $W(M) = \mathrm{Ker}\bigl(1 - p \colon \mathbf{G}_{a}^{n}
\to \mathbf{G}_{a}^{n}\bigr)$ est représentable.\footnote{La page note
$E^{n}$, sans le définir ; c'est l'espace affine $\mathbb{A}^{n}$, ou
$\mathbf{G}_{a}^{n}$ avec sa structure de module.}

En revanche $V(M)$ est représentable pour tout $M$, par l'algèbre symétrique
$\mathrm{Sym}^{*}_{k}(M) = \bigoplus_{n \geqslant 0} \mathrm{Sym}^{n}_{k}(M)$,
puisque $\mathrm{Hom}_{k}(M, k') = \mathrm{Hom}_{k\text{-alg}}(\mathrm{Sym}^{*}_{k} M, k')$.
L'accouplement $V(M) \times W(M) \to \mathcal{O} = W(k)$ donne un isomorphisme
\[
  V(M) \xrightarrow{\ \sim\ } \underline{\mathrm{Hom}}_{\mathcal{O}}\bigl(W(M), \mathcal{O}\bigr),
\]
parce que $M \mapsto W(M)$ est pleinement fidèle, et le reste après toute
extension de la base ; symétriquement, $V$ étant pleinement fidèle,
$W(M) \simeq \underline{\mathrm{Hom}}_{\mathcal{O}}(V(M), \mathcal{O})$. On a
donc une dualité parfaite entre $V(M)$ et $W(M)$, sans hypothèse de finitude
sur $M$.\footnote{La page conclut « Donc on a une dualité
parfaite entre $V(M)$ et $W(M)$ » ; le mot est incertain (peut-être
« bi-dualité »), et le second isomorphisme, qui résulte de la pleine fidélité
de $V$ (page 25), n'est pas écrit.}
Enfin $\underline{\mathrm{Aut}}_{\mathcal{O}}(M) \subset
\underline{\mathrm{End}}_{\mathcal{O}}(M)$ est représentable si $M$ est
projectif de type fini. La page s'arrête sur une ligne isolée, où se lisent
un mot incertain et $D(\Gamma)$ : sans doute l'annonce des groupes
diagonalisables.

\end{document}
